\summarychapter{27}{The Curse upon the Crown}{The Church of England}

\begin{inbrief}
The Lutheran reformation began in a monk's cell, the Anabaptist in a prayer meeting and the Calvinist
at a scholar's desk; the English reformation began in affairs of state, with the succession to the
throne. Under Henry VIII England rejected Rome without changing doctrine. The chapter follows the
break, the English Bible, the swings under Edward VI and Mary, and Elizabeth's middle way, the
background of Episcopalians, Congregationalists and Baptists.
\question{Why did England, without any great theological issue, overthrow the long-standing authority
of the Church of Rome?}
\end{inbrief}

On a Saturday afternoon in June 1533 Anne Boleyn rode through London in a chair draped in silver and
gold, on her way to be crowned queen at Westminster Abbey. Behind the procession lay the events that
carried England into the Reformation in its own way. England in a sense had two reformations, a
constitutional one under Henry VIII (1509--1547) and a theological one under the Puritans almost a
century later. In the first England forecast the future of Christianity in modern nations: belief
became almost entirely a private affair, and the practice of religion an instrument of the state,
what later generations called civil religion.

\ssection{England's Break with Rome}
The break grew out of the marital problems of Henry, whom Dickens called ``a blot of blood and grease
on the history of England.'' For centuries the church in England had been moving toward independence,
and patriotic Englishmen felt their faith to be distinctive. Cardinal Thomas Wolsey (1474--1530)
symbolizes this: archbishop of York, cardinal of Rome and chancellor of the realm, yet in every office
the king's man. The headline is that Henry revolted against the pope because he desired Anne Boleyn, a
lady of the court. The substance is the succession. Henry had mistresses and an illegitimate son;
what he lacked was a son by his queen, Catherine of Aragon, daughter of Ferdinand and Isabella. Of
her five children only the princess Mary survived, and England, whose one earlier queen had brought
wars of succession, would not accept a girl. By 1525 Catherine was forty, and Henry asked whether he
lay under God's curse, for she had briefly been the wife of his dead brother Arthur, and Leviticus
says that a man who takes his brother's wife shall be childless. Pope Julius II had given a special
permission for the marriage.

In 1527 Henry asked Clement VII to revoke that dispensation and declare the eighteen-year marriage
void. But Catherine was the aunt of the emperor Charles V, whom the pope could not afford to offend,
and he stalled. At the suggestion of Thomas Cranmer, Henry put his case to the universities of Europe;
their answers were mixed but gave him cover. In January 1533 he secretly married Anne, in May an
English church court annulled the first marriage, and in September Anne bore a child---a girl,
Elizabeth. The pope excommunicated him. Henry knew that antipapal feeling ran high (the Cambridge inn
where Luther's admirers met was called ``Little Germany'') and judged that he would meet little
resistance if he renounced the pope and left doctrine alone. Using a fourteenth-century law against
dealings with foreign powers, he forced the clergy to stop dealing with Rome. In 1534 the Act of
Supremacy declared the king ``the only supreme head in earth of the Church of England.'' England now
had a national church. Its head was no priest. He could appoint bishops but not consecrate them,
defend the faith but not formulate it, and for the priestly functions he looked to the archbishop of
Canterbury, Thomas Cranmer.

\ssection{Henry's Dual Policy}
The only religious issue was papal supremacy. Henry meant to keep the old faith and thought himself
its guardian. In 1521 he had answered Luther with a \emph{Defense of the Seven Sacraments}, calling
him a ``poisonous serpent,'' and the pope had named him ``Defender of the Faith,'' a title English
monarchs still bear. What he wanted was an English Catholic church in place of a Roman one. The
Statute of Six Articles (1539) upheld clerical celibacy, the private mass and confession to a priest.

Two things did change. The first was the suppression of the monasteries. Monks were neither popular
nor devout, and when the smaller houses were closed in 1536 fully half chose secular life over another
cloister. The monastic lands, almost a tenth of the nation's wealth, refilled the treasury and were
given or sold to barons and gentry; the monks received pensions. At one stroke Henry reduced
opposition and made friends. The second was the English Bible in the churches. Rome did not object in
principle to authorized translations, but versions made from Hebrew and Greek were bound to offend,
since important doctrines rested on doubtful renderings in the Latin, such as ``do penance'' where the
Greek says ``repent.'' After Erasmus's Greek New Testament of 1516, with its preface urging
translation, versions in German, French and English fed national feeling and Protestant conviction.

\ssection{William Tyndale's Bible}
The pioneer was William Tyndale. Amazed at the ignorance of the clergy, he told a fellow priest that
if God spared him he would make ``a boy behind the plow'' know more Scripture than he did. Unwelcome
in England, he fled to the Continent after study at Oxford and Cambridge, printed his New Testament
there, and early in 1526 began smuggling copies home. He went on to parts of the Old Testament and an
improved New. In 1536 he was caught, and after seventeen months in prison he was burned, praying,
``Lord, open the king of England's eyes.''

The answer was already coming. While he was in prison Miles Coverdale published the first complete
English Bible, basically Tyndale's work filled out from Latin and German versions. A year after his
death the Matthew Bible appeared, a careful compilation of Tyndale and Coverdale by John Rogers, who
left his own name off. At Cranmer's request Henry authorized this version, revised by Coverdale and
called the Great Bible, to be bought and read throughout the realm. The excitement was such that
Henry soon limited Bible reading to wealthy merchants and aristocrats, but wider freedom was only a
matter of time.

\ssection{The Swing to Protestantism}
Henry died in 1547 and was followed by his only son, the frail ten-year-old Edward VI, child of Jane
Seymour, whom Henry had married after executing Anne Boleyn for adultery. The advisers who governed
for Edward favored the Reformation, and policy swung sharply. The Six Articles were repealed, priests
were allowed to marry, and the Latin service gave way to Cranmer's \emph{Book of Common Prayer} in
English. In 1553 Cranmer's Forty-two Articles defined the church's faith on Protestant lines.

Edward died that year, and Catherine's daughter Mary, a devout Catholic, tried to lead England back
to Rome. In four years she outdid her father in intolerance and sent nearly three hundred Protestants
to the stake, Cranmer among them. John Foxe's \emph{Book of Martyrs} (1571) fixed a lasting horror of
Catholicism in the English mind and gave her the name Bloody Mary. Probably the only truly pious
English monarch of the century, she blundered in making martyrs. After a generation free of Rome the
country took her devotion and her marriage to Philip of Spain as betrayal, and she died broken.

Under Anne Boleyn's daughter Elizabeth I (1558--1603) the Anglican church took on its distinctive
character, neither Roman nor Reformed. Needing religious peace, she worked for a compromise. The
church stayed a state church under the monarch, but she changed her title from Supreme Head to the
more modest Supreme Governor. Her Thirty-nine Articles (1563) were essentially Protestant, taking the
Bible as final authority and recognizing only baptism and the Eucharist as sacraments instituted by
Christ, yet many were worded to satisfy both sides. The liturgy kept many Catholic elements, and
bishops in apostolic succession governed, for Elizabeth, like later monarchs, believed a church led
by bishops easier to control. Churchmen came to call this the \emph{Via Media}, the middle way. Some
of the exiles who returned from the Continent after Mary's reign were not convinced. They had read
their Bibles, had their own idea of a true reformation, and cried out against ``the ease in Zion.''
These were the Puritans.

\begin{keyfacts}
\keypeople{Henry VIII; Anne Boleyn; Catherine of Aragon; Cardinal Wolsey; Clement VII; Thomas Cranmer;
William Tyndale; Miles Coverdale; John Rogers; Edward VI; Mary I; John Foxe; Elizabeth I.}
\keydates{%
  \kd{1521}{Henry VIII named ``Defender of the Faith'' for writing against Luther}
  \kd{1526}{Tyndale's English New Testament smuggled into England}
  \kd{1533}{Henry VIII marries Anne Boleyn; his first marriage annulled}
  \kd{1534}{Act of Supremacy: the king supreme head of the Church of England}
  \kd{1536}{Smaller English monasteries suppressed; Tyndale burned}
  \kd{1539}{Statute of Six Articles upholds Catholic doctrine in England}
  \kd{1547}{Death of Henry VIII; Edward VI turns England toward Protestantism}
  \kd{1553}{Forty-two Articles; Mary I restores Catholicism in England}
  \kd{1558--1603}{Reign of Elizabeth I}
  \kd{1563}{Thirty-nine Articles of the Church of England}}
\keyterms{Act of Supremacy; civil religion; Defender of the Faith; dissolution of the monasteries;
Great Bible; \emph{Book of Common Prayer}; Thirty-nine Articles; Supreme Governor; \emph{Via Media};
Puritans.}
\end{keyfacts}
